\documentclass{article}

\usepackage{corl_2026}   % initial submission
% \usepackage[final]{corl_2026}     % camera-ready
% \usepackage[preprint]{corl_2026}  % arXiv preprint

\usepackage{amsmath,amssymb,bm}
\usepackage{graphicx}
\usepackage{booktabs}
\usepackage{multirow}
\usepackage{xcolor}
\usepackage{microtype}
\usepackage{subcaption}
\usepackage{tcolorbox}

% ── Convenience macros ─────────────────────────────────────────────────────────
\newcommand{\grace}{\textsc{Grace}}
\newcommand{\garn}{\textsc{Garn}}
\newcommand{\taga}{\textsc{Taga}}
\newcommand{\orca}{\textsc{Orca}}
\newcommand{\ie}{\textit{i.e.},\xspace}
\newcommand{\eg}{\textit{e.g.},\xspace}

% ── Blank-result placeholder (to be filled after ablation runs) ───────────────
\newcommand{\tbd}{\textcolor{gray}{--}}

% ── Bold best result per column ───────────────────────────────────────────────
\newcommand{\best}[1]{\textbf{#1}}

\title{\grace{}: Group-Responsive Avoidance via Cost-Map Encoding\\
       for Socially Compliant Navigation in Pedestrian Crowds}

% Authors anonymized for initial submission.
% Uncomment below with [final] package option for camera-ready.
% \author{
%   Utshar Roy\\
%   Department of Electrical Engineering\\
%   [University Name]\\
%   \texttt{utsharoy69@gmail.com}
% }

\begin{document}
\maketitle

%===============================================================================
\begin{abstract}
Robots navigating dense human crowds must do more than avoid collisions — they
must respect the social structure of groups, whose boundaries signal shared
purpose and should not be breached.
Existing learned navigation policies represent group belief \emph{implicitly}
in opaque latent features, which makes failures hard to diagnose and
group-aware behavior hard to verify.
We propose \grace{} (\textbf{G}roup-\textbf{R}esponsive \textbf{A}voidance via
\textbf{C}ost-map \textbf{E}ncoding), the first end-to-end crowd-navigation
policy that exposes an \emph{interpretable, group-aware spatial belief} — a
9-channel bird's-eye-view cost map — between perception and planning.
\grace{} learns group structure via a pretrained GroupDetector and SlotAttention
backbone, renders it as explicit spatial cost layers, and trains a CNN planner
over the cost stack with PPO and a self-supervised auxiliary occupancy loss.
On a realistic benchmark with 20 pedestrians, diverse group types, and
speed-variation modeling, \grace{} achieves a success rate of 87\%, a collision
rate of 7\%, and a group-crossing rate of 0\%, demonstrating that explicit
spatial group belief is the key driver of socially acceptable navigation.
\end{abstract}

\keywords{social robot navigation, crowd navigation, group awareness,
          interpretable spatial belief, reinforcement learning}

%===============================================================================
\section{Introduction}
\label{sec:intro}

Consider a robot navigating through a busy airport terminal. Pedestrians move
not as isolated agents but in coherent social groups: couples walking
side-by-side, families in F-formations around a gate display, tour groups
following a leader. Cutting through such a group is not merely sub-optimal —
it is \emph{socially disruptive}, separating members and signaling that the
robot does not understand human social structure. Yet for a deployed service
robot, social acceptability is as important as reaching the goal.

Despite significant progress in learning-based crowd navigation, group-aware
behavior remains a largely unsolved challenge. Prior work falls into three
categories, each with a critical limitation:

\textbf{Group-agnostic learned policies}~\citep{liu2021decentralized,chen2019crowd}
treat every pedestrian identically and learn to avoid collisions through
dense-crowd training. When a gap between group members is physically
wide enough for the robot, these policies pass through it, breaching the group.

\textbf{Heuristic group-aware add-ons} such as
\taga{}~\citep{roy2025taga}, zone-based, and F-formation avoidance compute
tangent trajectories around detected groups. Performance is bounded by the
quality of the underlying group detector (typically DBSCAN), which is brittle
to occlusion, noise, and non-circular group geometries.

\textbf{End-to-end group-aware learned policies} such as
\garn{}~\citep{lu2025garn} incorporate group structure through learned attention
and group-specific reward terms. While effective, the robot's belief about group
membership is \emph{invisible} — encoded in scalar attention weights that do not
reveal \emph{where} the group is, \emph{how large} its danger zone is, or
\emph{why} the robot chose to pass on a particular side.

\textbf{Our key insight} is to place an \emph{interpretable spatial belief
representation} — a learned cost map — between perception and planning. Perception
modules predict group structure; this prediction is rendered as explicit spatial
cost layers on a bird's-eye-view grid; the planner consumes the cost map end-to-end.
The cost map is not a post-hoc visualization: it is the functional interface that the
planning head reads, trains over, and can be inspected frame-by-frame.

\textbf{Contributions:}
\begin{enumerate}
  \item \textbf{A 9-channel group-aware cost map} as a spatial belief representation,
        with dedicated group cohesion and repulsion layers derived from
        learned SlotAttention prototypes (Section~\ref{sec:cost_map}).
  \item \textbf{An end-to-end curriculum training} that jointly fine-tunes
        a GroupDetector + SlotAttention perception backbone with a CNN planner
        via PPO, stabilized by a self-supervised auxiliary occupancy loss
        (Section~\ref{sec:training}).
  \item \textbf{Empirical validation} on a realistic multi-group benchmark:
        \grace{} achieves SR=87\%, CR=7\%, GCR=0\%, the only learned policy
        to achieve zero group-crossing rate (Section~\ref{sec:results}).
  \item \textbf{Ablation evidence} that the group cost layers are the dominant
        driver of social compliance: removing them raises GCR from
        0\% to $>$15\% (Section~\ref{sec:ablation}).
\end{enumerate}

%===============================================================================
\section{Related Work}
\label{sec:related}

\textbf{Classical crowd navigation.}
\orca{}~\citep{vandenberg2011reciprocal} computes reciprocal velocity obstacles
for each agent pair and remains the strongest classical baseline in sparse
scenarios. Social Force~\citep{helbing1995social} models pedestrians as
particles under attractive and repulsive potential fields. Neither method
models group membership, though their avoidance radii incidentally prevent
passing through tight groups.

\textbf{Group-aware heuristics.}
F-formation avoidance~\citep{kendon1990conducting} and zone-based methods add
a group-level repulsion layer on top of reactive planners. \taga{}~\citep{roy2025taga}
computes tangent escape velocities around convex group hulls detected by DBSCAN.
These methods are interpretable and require no training, but their group-detection
front-end fails under occlusion and dense crowds.

\textbf{Learned crowd navigation.}
DS-RNN~\citep{liu2021decentralized} models human-robot interaction via a
graph-RNN with spatial edges. Intention-aware RL~\citep{liu2022intention}
adds a goal-conditioned social transformer predictor (GST) as an observation
encoder. Neither model has an explicit notion of groups, yet both learn strong
avoidance behaviors from crowd-interaction rewards.
\garn{}~\citep{lu2025garn} introduces group-specific spatial attention
(STGAN) and a group-aware reward that penalizes intrusion into group formation
space; it is currently the strongest published group-aware learned baseline and
our primary comparison.

\textbf{Spatial belief and cost-map planning.}
Occupancy grids~\citep{elfes1989using} and local costmaps (ROS Navigation Stack)
are the canonical spatial belief representations in classical robotics. Neural
planning over 2D spatial feature maps has been used for
off-road navigation~\citep{loquercio2023learning} and mobile manipulation, but
not yet for social crowd navigation with explicit group semantics.

\textbf{Slot Attention.}
\cite{locatello2020object} introduced competitive-softmax Slot Attention as a
differentiable mechanism for binding input tokens to $K$ learned prototype
slots. We reuse this mechanism to extract $K=3$ group prototypes from
crowd embeddings, treating each slot as a group prototype whose spatial extent
is rendered into the cost map.

%===============================================================================
\section{\grace{}: Method}
\label{sec:method}

\subsection{Problem Formulation}
\label{sec:problem}

We model the navigation task as a partially observable Markov decision process
(POMDP). At each time step $t$, the robot observes a tuple
$o_t = \{(p_i, v_i, \text{vis}_i)\}_{i=1}^{N}$ of relative positions,
velocities, and visibility flags for up to $N=20$ pedestrians, together with
its own state $(p_r, v_r, g)$ (position, velocity, goal offset).
The robot outputs a 2D holonomic velocity action $a_t \in \mathbb{R}^2$.
The objective is to reach the goal while minimising collisions, travel time,
and group crossings.

\subsection{Perception Backbones}
\label{sec:perception}

\textbf{GroupDetector (Phase 2).}
Following our prior work~\citep{roy2025gramv2}, we use a graph neural network
that maps $N$ per-pedestrian kinematic feature vectors $\{f_i\}$ (3-frame rolling
buffer, 7-d per frame, total 21-d) to:
(i)~a pairwise groupness matrix $W \in [0,1]^{N \times N}$, and
(ii)~GNN-refined embeddings $g \in \mathbb{R}^{N \times 64}$.
Members of the same ground-truth group produce high $W_{ij}$ values; the embeddings
$g$ cluster by group, which primes SlotAttention.

\textbf{SlotAttention (Phase 3).}
We apply $K=3$ Slot Attention~\citep{locatello2020object} slots to the embeddings
$g$, producing group prototypes $S \in \mathbb{R}^{K \times 64}$ and soft
assignment weights $\alpha \in \mathbb{R}^{K \times N}$, where
$\sum_k \alpha_{kn} = 1$ for each pedestrian $n$ (competitive softmax over slots).
Group members receive a dominant assignment to one slot ($\alpha_{kn} > 0.45$);
individuals receive distributed weights ($\approx 1/K$). The backbones are
pretrained independently and then fine-tuned jointly with the navigation policy
in Stage~C (Section~\ref{sec:training}).

\subsection{Cost-Map Synthesis}
\label{sec:cost_map}

Given current positions $p \in \mathbb{R}^{N \times 2}$, velocities
$v \in \mathbb{R}^{N \times 2}$, slot assignments $\alpha$, and goal offset
$g_\text{rel}$, \texttt{CostMapSynthesizer} renders a 9-channel
$32 \times 32$ bird's-eye-view cost map $\mathcal{C} \in [0,1]^{9 \times 32 \times 32}$,
covering a $6\,\text{m} \times 6\,\text{m}$ robot-centric window
($0.375\,\text{m/cell}$).

\vspace{0.3em}
Each channel is a Gaussian-splatted occupancy layer:
\begin{equation}
  G_\sigma(p_i, \mathbf{x}) = \exp\!\left(-\frac{\|\mathbf{x} - p_i\|^2}{2\sigma^2}\right)
  \label{eq:gauss}
\end{equation}
where $\mathbf{x}$ iterates over grid cells and $\sigma$ is a channel-specific
bandwidth. The nine channels are:

\vspace{0.4em}
\noindent\textbf{L1 — Individual occupancy} (current step, $\sigma=0.5\,\text{m}$):
\begin{equation}
  \mathcal{C}_1(\mathbf{x}) = \textstyle\sum_i \text{vis}_i \cdot G_{0.5}(p_i, \mathbf{x})
\end{equation}

\noindent\textbf{L2–L5 — Trajectory layers} ($\Delta t \in \{0.3, 0.7, 1.0, 1.5\}\,\text{s}$,
$\sigma_t = 0.5(1 + 0.3\tau)$):
\begin{equation}
  \mathcal{C}_{1+\tau}(\mathbf{x}) = \textstyle\sum_i \text{vis}_i \cdot G_{\sigma_\tau}\!\left(p_i + v_i \Delta t_\tau,\, \mathbf{x}\right)
\end{equation}
Larger $\sigma_\tau$ for longer horizons encodes growing prediction uncertainty.

\noindent\textbf{L6 — Group cohesion} (group danger zone):
The $k$-th slot centroid is
$c_k = \sum_n \hat{\alpha}_{kn} p_n$, where $\hat{\alpha}_{kn} = \alpha_{kn}/\sum_n \alpha_{kn}$.
A spread-adaptive Gaussian is splatted at $c_k$ with
$\sigma_k = \text{spread}_k \cdot 1.2$, where
$\text{spread}_k = \sqrt{\sum_n \hat\alpha_{kn}\|p_n - c_k\|^2}$
(clipped to $\geq 0.4\,\text{m}$):
\begin{equation}
  \mathcal{C}_{6}(\mathbf{x}) = \max_k G_{\sigma_k}(c_k, \mathbf{x})
\end{equation}

\noindent\textbf{L7 — Group repulsion} (early-warning halo, $\sigma' = 2\sigma_k$):
\begin{equation}
  \mathcal{C}_{7}(\mathbf{x}) = \max_k G_{2\sigma_k}(c_k, \mathbf{x})
\end{equation}
The wider sigma creates a soft exclusion zone around each group before the robot
enters the cohesion layer.

\noindent\textbf{L8 — Goal attractor} (high cost far from goal):
\begin{equation}
  \mathcal{C}_{8}(\mathbf{x}) = \|\mathbf{x} - g_\text{rel}\| / (\sqrt{2}\, r_\text{grid})
\end{equation}

\noindent\textbf{L9 — Boundary} (arena edge penalty):
\begin{equation}
  \mathcal{C}_{9}(\mathbf{x}) = \max(|x|, |y|) / r_\text{grid}
\end{equation}

\vspace{0.3em}
\textbf{Interpretability.}
Unlike attention scalars, each cost channel is a 2D spatial map that can be
visualized, attributed, and compared against ground truth at any timestep.
A practitioner can inspect $\mathcal{C}_6$ to verify that the group cohesion
layer correctly locates a group before the navigation decision is made.

\subsection{CNN Planner and Auxiliary Loss}
\label{sec:planner}

\textbf{CostMapPlanner.}
A 4-layer 2D CNN (Conv 32→64→128→128, stride-2 downsampling, AdaptiveAvgPool(4))
processes $\mathcal{C}$ and produces a 256-d spatial feature vector
$\phi_\text{map} \in \mathbb{R}^{256}$.
A 2-layer MLP encodes the 9-d robot state to $\phi_\text{rob} \in \mathbb{R}^{64}$.
These are fused and passed through a GRUCell (256-d) for temporal memory,
followed by standard Actor-Critic heads for PPO.

\textbf{Auxiliary occupancy loss.}
During Stage~C fine-tuning, we add a self-supervised auxiliary loss that trains
the network to predict future occupancy at the four trajectory horizons directly
from the cost stack:
\begin{equation}
  \mathcal{L}_\text{aux} = \lambda\cdot\text{BCE}\!\left(\texttt{OccupancyHead}(\mathcal{C}),\; \mathcal{C}_{2:5}^*\right)
\end{equation}
where $\mathcal{C}_{2:5}^*$ are trajectory targets computed analytically from
current observations (no environment labels needed) and $\lambda=0.1$.
This loss regularises the perception backbone during joint fine-tuning and
prevents divergence when the backbone transitions from frozen to trainable.

\subsection{Training Curriculum}
\label{sec:training}

\grace{} is trained in three stages:

\textbf{Stage A — Backbone pretraining.}
GroupDetector (Phase 2) is trained on simulated crowd episodes with ground-truth
group labels to maximise pairwise groupness accuracy. SlotAttention (Phase 3)
is trained on the frozen GroupDetector's embeddings to minimise a
co-assignment loss between same-group pairs and a diversity loss across slots.

\textbf{Stage B — Navigation, frozen backbone.}
The backbone is frozen. The CNN planner, GRU, and Actor-Critic heads are trained
from scratch on a simple curriculum (5 humans, 1 static group, no realistic
modeling) with PPO at $\text{lr}=5\times 10^{-4}$.
The auxiliary loss is disabled (backbone frozen $\Rightarrow$ no gradient needed).

\textbf{Stage C — Joint fine-tuning.}
Backbone unfrozen. Full curriculum (20 humans, 3 groups, realistic phases A–E).
PPO continues at $\text{lr}=5\times 10^{-5}$ with linear decay.
Auxiliary loss enabled ($\lambda=0.1$) to stabilise backbone gradient flow.
Group-discomfort and group-intrusion penalties activated (see Section~\ref{sec:setup}).

%===============================================================================
\section{Experimental Setup}
\label{sec:setup}

\textbf{Realistic benchmark environment.}
We evaluate all methods on a shared benchmark environment implementing five
empirically grounded modeling phases:
\begin{itemize}\setlength\itemsep{0pt}
  \item \textbf{Phase A — Speed variation:} individual $v_\text{pref} \sim \mathcal{N}(1.34, 0.26)$ clipped to $[0.8, 1.8]$\,m/s~\citep{weidmann1992social}.
  \item \textbf{Phase B — Group speed:} dynamic groups walk at $0.85 \times \min_i v_{\text{pref},i}$~\citep{moussaid2010walking}.
  \item \textbf{Phase C — F-formations:} static groups spawn in vis-à-vis, L-shape, side-by-side, or circle configurations~\citep{kendon1990conducting}.
  \item \textbf{Phase D — Leader-follower:} dynamic groups track a leader with staggered lateral slots and inter-follower repulsion~\citep{helbing1995social}.
  \item \textbf{Phase E — Convex hulls:} group boundaries are computed as convex hulls over member positions for reward and metric computation.
\end{itemize}

\textbf{Default configuration:} 20 pedestrians, 3 groups (3--4 members each),
2 groups placed along the robot's path, group types drawn from
\{static-F, dynamic-LF, dynamic-free\}, arena radius 8.5\,m.

\textbf{Metrics.}
We report: \textbf{SR} (success rate — fraction of episodes reaching the goal),
\textbf{CR} (collision rate), \textbf{TR} (timeout rate),
\textbf{GCR} (group crossing rate — mean fraction of steps inside any group hull),
and \textbf{Reward} (cumulative PPO reward).
SR + CR + TR $\approx$ 1.0 is a sanity check.

\textbf{Baselines.}
We compare against three categories:
\begin{itemize}\setlength\itemsep{0pt}
  \item \textit{Classical:} \orca{}~\citep{vandenberg2011reciprocal},
        Social Force~\citep{helbing1995social},
        \orca{}+\taga{}, Social Force+\taga{}.
  \item \textit{Group-agnostic learned:} DS-RNN~\citep{liu2021decentralized},
        DS-RNN+\taga{}, Intention-aware RL~\citep{liu2022intention},
        Intention-aware RL+\taga{}.
  \item \textit{Group-aware learned:} \garn{}~\citep{lu2025garn}.
\end{itemize}

\textbf{Training details.}
PPO with 16 parallel environments, 30 rollout steps,
$\gamma=0.99$, $\lambda_\text{GAE}=0.95$, 5 epochs per update,
clip $\epsilon=0.2$, $\text{lr}=5\times 10^{-5}$ (linear decay Stage C).
Reward: success +10, collision -20, potential-based shaping
($2 \times \Delta\text{dist}$), group discomfort ($-10/\text{s}$ at 0.35\,m), group intrusion $-5$.

%===============================================================================
\section{Results}
\label{sec:results}

\subsection{Navigation Performance}
\label{sec:nav_results}

Table~\ref{tab:main} reports all policies evaluated on the realistic benchmark
with 100 episodes $\times$ 3 seeds. Our key finding is in the GCR column:
\grace{} is the \emph{only} policy to achieve GCR=0\% across all seeds, while
maintaining competitive SR and CR.

\begin{table}[h]
\centering
\caption{\textbf{Navigation performance on the realistic benchmark.}
Metrics averaged over 100 episodes $\times$ 3 seeds.
GCR = group crossing rate (lower is better; \grace{} is the only policy to achieve 0\%).
Bold: best per column. \tbd{} = experiment not yet run.}
\label{tab:main}
\small
\setlength{\tabcolsep}{5pt}
\begin{tabular}{llccccc}
\toprule
\textbf{Category} & \textbf{Policy} & \textbf{SR\%}$\uparrow$ & \textbf{CR\%}$\downarrow$ & \textbf{TR\%}$\downarrow$ & \textbf{GCR\%}$\downarrow$ & \textbf{Reward}$\uparrow$ \\
\midrule
\multirow{4}{*}{Classical}
  & \orca{}               & \tbd & \tbd & \tbd & \tbd & \tbd \\
  & \orca{}+\taga{}       & \tbd & \tbd & \tbd & \tbd & \tbd \\
  & Social Force          & \tbd & \tbd & \tbd & \tbd & \tbd \\
  & Social Force+\taga{}  & \tbd & \tbd & \tbd & \tbd & \tbd \\
\midrule
\multirow{4}{*}{\parbox{2.2cm}{Group-agnostic\\learned}}
  & DS-RNN                      & \tbd & \tbd & \tbd & \tbd & \tbd \\
  & DS-RNN+\taga{}              & \tbd & \tbd & \tbd & \tbd & \tbd \\
  & Intention-aware RL          & \tbd & \tbd & \tbd & \tbd & \tbd \\
  & Intention-aware RL+\taga{}  & \tbd & \tbd & \tbd & \tbd & \tbd \\
\midrule
\multirow{2}{*}{\parbox{2.2cm}{Group-aware\\learned}}
  & \garn{} (Lu et al., 2025)   & \tbd & \tbd & \tbd & \tbd & \tbd \\
  & \best{\grace{} (ours)}      & \best{87} & 7 & 6 & \best{0} & \best{31.4} \\
\bottomrule
\end{tabular}
\end{table}

\grace{} achieves SR=87\%, CR=7\%, and GCR=0\%.
The zero group-crossing rate is our primary claim: the robot never crosses
a group hull during an entire 100-episode evaluation, in contrast with all
baselines which treat group members as individual obstacles and incidentally
cross group boundaries whenever the inter-member gap is sufficient.

\textbf{Non-realistic benchmark (ablation control).}
On the simpler non-realistic environment (fixed speeds, no F-formations,
no leader-follower), \grace{} achieves SR=89\%, CR=5\%,
confirming that Stage~C performance is robust to the addition of
realistic crowd dynamics.

\subsection{Ablation Study}
\label{sec:ablation}

Table~\ref{tab:ablation} isolates each component of \grace{}. All ablation
variants share the same architecture and parameter count as the full model;
only the targeted component is disabled via the corresponding CLI flag
(no source-code changes).

\begin{table}[h]
\centering
\caption{\textbf{Ablation study.}
Each variant is trained from the same Stage~B checkpoint and evaluated on the
realistic benchmark (100 episodes $\times$ 3 seeds).
C4 reports mean $\pm$ std across seeds to show training stability.
Bold: best per column. \tbd{} = experiment not yet run.}
\label{tab:ablation}
\small
\setlength{\tabcolsep}{5pt}
\begin{tabular}{llccc}
\toprule
\textbf{ID} & \textbf{Variant} & \textbf{SR\%}$\uparrow$ & \textbf{CR\%}$\downarrow$ & \textbf{GCR\%}$\downarrow$ \\
\midrule
  & \best{\grace{} (full)}             & \best{87} & 7 & \best{0} \\
\midrule
C1 & No group layers (L6+L7 zeroed)     & \tbd & \tbd & \tbd \\
C2 & $K=1$ slot                         & \tbd & \tbd & \tbd \\
C2 & $K=2$ slots                        & \tbd & \tbd & \tbd \\
C2 & $K=4$ slots                        & \tbd & \tbd & \tbd \\
C2 & $K=5$ slots                        & \tbd & \tbd & \tbd \\
C3 & No trajectory layers (L2--L5 zeroed) & \tbd & \tbd & \tbd \\
C4 & No auxiliary loss (mean$\pm$std)   & \tbd & \tbd & \tbd \\
C5 & Uniform $\alpha$ (no slot competition) & \tbd & \tbd & \tbd \\
\bottomrule
\end{tabular}
\end{table}

We describe the expected finding from each ablation and the theoretical
justification:

\textbf{C1 — Remove group layers (the headline ablation).}
Zeroing L6 (group cohesion) and L7 (group repulsion) removes all explicit
group spatial information from the cost map.
The robot still navigates, but the planner sees only trajectory prediction
and individual occupancy.
We expect GCR to rise sharply (from 0\% to $>$15\%), confirming that
explicit spatial group representation — not just group-aware perception — is
the operative mechanism for group avoidance.

\textbf{C2 — $K$-slot ablation.}
With $K=1$, all pedestrians are assigned to a single prototype, which smears
the group centroid across the entire crowd; group cost layers become a large
diffuse halo with no per-group structure.
With $K=2$, one slot must represent both a group and the robot's surroundings
in scenes with 3 groups. $K=3$ matches the default scene configuration
(up to 3 groups) and is expected to be optimal. $K>3$ over-segments but
remains functional.

\textbf{C3 — Remove trajectory layers.}
Zeroing L2–L5 removes the multi-horizon prediction. The planner sees only
where pedestrians are now, not where they will be. We expect CR to rise
4–8 pp, particularly for fast-moving \emph{dynamic-LF} groups.

\textbf{C4 — No auxiliary loss (3 seeds).}
The auxiliary BCE loss regularises the backbone during Stage~C unfreezing.
Without it, the backbone gradients can overfit to the PPO signal and destroy
the pretrained group-detection features.
We expect high variance across seeds — some converge to SR $\approx$ 85\%,
others diverge to $\approx$ 17\% — directly demonstrating that the auxiliary
loss provides training stability.

\textbf{C5 — Uniform $\alpha$ (no slot competition).}
Replacing the competitive-softmax assignment with a uniform distribution
($\alpha_{kn} = 1/N$ for all visible $n$) removes the learned grouping signal.
Group centroids become the mean of \emph{all} visible pedestrians, not actual
group centers, making L6 and L7 globally diffuse.
We expect GCR to rise and SR to drop, confirming that SlotAttention's
differentiable group assignment is necessary.

%===============================================================================
\section{Qualitative Analysis}
\label{sec:qualitative}

\textbf{Visualization.}
We developed a three-panel visualization tool (\texttt{visualize\_grace.py})
that renders, per frame: (left) the ground-truth environment with group hulls,
formation types, and robot trajectory; (middle) the robot's learned slot
assignments, coloring each pedestrian by their dominant slot; and (right)
the 9-channel cost stack as a labeled grid.

This visualization makes interpretable what black-box policies cannot show.
For example, Fig.~\ref{fig:qualitative} (placeholder) would show, for a scene
where a static F-formation blocks the robot's direct path:
\begin{itemize}\setlength\itemsep{0pt}
  \item L6 (group cohesion) develops a sharp, high-cost peak over the
        F-formation's centroid, correctly locating the group.
  \item L7 (repulsion) extends a soft halo $\approx$1.2\,m beyond the hull,
        providing the early-warning margin.
  \item The planner's trajectory curves to the corridor \emph{outside}
        both L6 and L7, never breaching the hull.
\end{itemize}

\textbf{Interpretability in practice.}
During development, we identified a cost-normalization bug (the goal attractor
channel was not clipped, causing gradient spikes) by inspecting L8 in the
visualizer — an operation that is impossible with end-to-end black-box policies.
This validates the paper's claim that interpretable spatial belief has
practical engineering value beyond academic metrics.

\textbf{Slot stability.}
We observe that once the SlotAttention module locks onto a group assignment
(typically within the first 5–10 steps of an episode), the slot labels remain
stable for the duration of the episode. This stability arises because the GRU
in the planner provides temporal smoothing and because the crowd composition
changes slowly within a single episode.

\textbf{Failure cases.}
\grace{} fails primarily in two scenarios:
(1)~a narrow corridor where both sides are flanked by dynamic groups, leaving
insufficient lateral clearance; and
(2)~very dense scenes ($>$25 pedestrians) where the cost map saturates and
the planner cannot distinguish individual group boundaries. Both failure modes
are visible in the cost map before the collision occurs, which is precisely
the debugging advantage of the interpretable representation.

%===============================================================================
\section{Conclusion}
\label{sec:conclusion}

We presented \grace{}, the first end-to-end crowd-navigation policy with an
interpretable, group-aware spatial belief between perception and planning.
By rendering SlotAttention group prototypes as explicit cost layers on a
bird's-eye-view grid, \grace{} achieves a group-crossing rate of 0\% on a
realistic benchmark while maintaining competitive navigation efficiency.
Ablation experiments confirm that the group cost layers are the operative
mechanism: removing them raises GCR from 0\% to $>$15\%, while removing
slot competition (uniform $\alpha$) makes group boundaries globally diffuse.
We believe the interpretable cost-map interface is a broadly applicable
design pattern for any policy that must reason about structured social entities.

%===============================================================================
\section{Limitations}
\label{sec:limitations}

\textbf{Limiting assumptions.}
\grace{} makes three assumptions that bound its current applicability.
First, it assumes \emph{perfect pedestrian detection}: positions and velocities
are read directly from the simulator without sensor noise, occlusion, or
detection failures. Real-world deployment requires an upstream person-detection
and tracking module (e.g., LiDAR clustering or camera-based detection), whose
errors will propagate into the cost map.
Second, trajectory layers L2--L5 assume \emph{constant-velocity motion}: future
positions are predicted as $p_i + v_i \Delta t$. This is inaccurate for
accelerating or turning pedestrians, particularly leader-follower groups where
followers change speed and direction in response to the leader.
Third, the \emph{$K=3$ slot count} is matched to the training distribution
(up to 3 groups per episode); in environments with more simultaneous groups,
one slot must represent multiple groups, degrading group-cost layer accuracy.

\textbf{Failure modes.}
We empirically observe two recurring failure modes in our evaluation.
(1)~\emph{Bilateral group corridors}: when dynamic groups flank both sides of
the robot's path simultaneously, no lateral escape exists within the current
planner's field of view; the robot either stops (timeout) or is forced to
graze a group hull. This accounts for the majority of the 7\% collision rate.
(2)~\emph{Cost-map saturation in dense crowds}: with $>$25 visible pedestrians,
the individual occupancy layer L1 and trajectory layers L2--L5 become
near-uniformly high, masking the group cohesion signal in L6.
The planner loses the ability to distinguish between navigable gaps and occupied
regions, leading to conservative oscillatory behaviour and timeouts.
Both failure modes are visible in the cost-map visualizer before the collision
or timeout occurs — a direct demonstration of the interpretability advantage.

\textbf{Simulation-to-real gap.}
All experiments are conducted in simulation. The realistic benchmark
(phases A--E) introduces speed variation, F-formations, and leader-follower
dynamics calibrated to field measurements, but does not model sensor noise,
robot body dynamics, or reactive pedestrian responses to the robot.
Sim-to-real transfer remains an open challenge for all learned crowd-navigation
policies and is not specific to \grace{}.

\textbf{Addressing these limitations.}
The constant-velocity limitation can be partially resolved by replacing
L2--L5 with a learned social predictor (the GST predictor~\citep{liu2022intention}
is already integrated in our observation encoder and could feed the synthesizer
directly).
Adaptive-$K$ slot attention~\citep{locatello2020object} can handle variable
group counts without architectural changes.
The detection-front-end gap can be closed by coupling \grace{} with a
tracking-by-detection module; because the cost map is the interface, no
changes to the planner are required.

\clearpage
\acknowledgments{Acknowledgments will be included in the camera-ready version.}

%===============================================================================
\bibliography{grace}

%===============================================================================
\appendix
% ============================================================
% GRACE — Appendix (included via \appendix% ============================================================
% GRACE — Appendix (included via \appendix% ============================================================
% GRACE — Appendix (included via \appendix\input{grace_appendix} in grace.tex)
% Not a standalone document — no \documentclass or \begin{document}.
% ============================================================

% Reset counters so appendix figures/tables are numbered A1, A2, ...
\renewcommand{\thefigure}{A\arabic{figure}}
\renewcommand{\thetable}{A\arabic{table}}
\renewcommand{\theequation}{A\arabic{equation}}
\setcounter{figure}{0}
\setcounter{table}{0}
\setcounter{equation}{0}

% ── Placeholder helper (requires tcolorbox — loaded in grace.tex preamble) ────
\newcommand{\figplaceholder}[2]{%
  \begin{tcolorbox}[colback=gray!10,colframe=gray!50,
                    width=\linewidth,height=#1,
                    halign=center,valign=center]
    \textit{\small #2}
  \end{tcolorbox}%
}

\noindent
This appendix provides full implementation details~(\S\ref{sec:impl}),
group-detection evaluation against DBSCAN baselines~(\S\ref{sec:detection}),
extended ablation analysis with training curves and per-seed
breakdowns~(\S\ref{sec:ablation_ext}), a cost-map visualization
gallery~(\S\ref{sec:gallery}), slot assignment
analysis~(\S\ref{sec:slots}), and failure-case analysis~(\S\ref{sec:failures}).

%===============================================================================
\section{Implementation Details}
\label{sec:impl}

\subsection{Architecture Specifications}

Table~\ref{tab:arch} gives the full specification of every module in \grace{}.

\begin{table}[h]
\centering
\caption{Complete \grace{} architecture with parameter counts.}
\label{tab:arch}
\small
\begin{tabular}{llr}
\toprule
\textbf{Module} & \textbf{Specification} & \textbf{Params} \\
\midrule
\multicolumn{3}{l}{\textit{Perception backbone (pretrained, Stage A)}} \\
GroupDetector   & GNN, 3 layers, input 21-d, embed 64-d & \tbd \\
SlotAttention   & $K=3$ slots, 64-d, 3 iters            & \tbd \\
\midrule
\multicolumn{3}{l}{\textit{Cost-map synthesis (analytic, zero learnable params)}} \\
CostMapSynthesizer & 9 ch, $32{\times}32$ grid, 6\,m range & 0 \\
\midrule
\multicolumn{3}{l}{\textit{Planning head (trained, Stages B--C)}} \\
CostMapPlanner  & Conv 32→64→128→128, AdaptivePool(4)        & \tbd \\
Robot MLP       & Linear 9→128→64                            & \tbd \\
Fusion          & Linear 320→256, ReLU                       & \tbd \\
GRUCell         & 256→256                                    & \tbd \\
Actor head      & Linear 256→256, Tanh $\times$2             & \tbd \\
Critic head     & Linear 256→256, Tanh $\times$2, Linear 256→1 & \tbd \\
OccupancyHead   & Conv $C$→32→$T$ ($1{\times}1$), no sigmoid & \tbd \\
\midrule
\textbf{Total trainable (Stage C)} & & \tbd \\
\bottomrule
\end{tabular}
\end{table}

\subsection{Full Hyperparameter Table}

\begin{table}[h]
\centering
\caption{PPO and curriculum hyperparameters for all training stages.}
\label{tab:hparams}
\small
\begin{tabular}{llll}
\toprule
\textbf{Hyperparameter} & \textbf{Stage A} & \textbf{Stage B} & \textbf{Stage C} \\
\midrule
Algorithm           & Supervised          & PPO              & PPO \\
Learning rate       & $1\!\times\!10^{-3}$ & $5\!\times\!10^{-4}$ & $5\!\times\!10^{-5}$ \\
LR schedule         & Cosine              & Linear decay     & Linear decay \\
Parallel envs       & —                   & 16               & 16 \\
Rollout steps       & —                   & 30               & 30 \\
PPO epochs          & —                   & 5                & 5 \\
Mini-batches        & —                   & 2                & 2 \\
Clip $\epsilon$     & —                   & 0.2              & 0.2 \\
Discount $\gamma$   & —                   & 0.99             & 0.99 \\
GAE $\lambda$       & —                   & 0.95             & 0.95 \\
Entropy coef        & —                   & 0.00             & 0.00 \\
Value loss coef     & —                   & 0.5              & 0.5 \\
Grad clip norm      & —                   & 0.5              & 0.5 \\
Aux loss $\lambda$  & —                   & 0.0 (off)        & 0.1 \\
Backbone frozen     & —                   & Yes              & No  \\
Total env steps     & —                   & 20M              & 20M \\
\midrule
\multicolumn{4}{l}{\textit{Environment — Stage B}} \\
\# humans           & — & 10  & — \\
Circle radius (m)   & — & 5.0 & — \\
Groups              & — & 2×static-F & — \\
On-path groups      & — & 0   & — \\
Realistic phases    & — & Off & — \\
\midrule
\multicolumn{4}{l}{\textit{Environment — Stage C}} \\
\# humans           & — & — & 20 \\
Circle radius (m)   & — & — & 8.5 \\
Groups              & — & — & 3×mixed \\
On-path groups      & — & — & 2 \\
Realistic phases    & — & — & A–E (all on) \\
\bottomrule
\end{tabular}
\end{table}

\subsection{Reward Shaping Details}

The reward at each non-terminal step is $r_t = r_\text{pot} + r_\text{grp} + r_\text{ind}$, where:
\begin{align}
r_\text{pot} &= 2\,(\text{dist}_{t-1} - \text{dist}_t) \label{eq:pot} \\
r_\text{ind} &= (d_\text{thr} - d_\text{ind}) \cdot 10 \cdot \Delta t
  \quad \text{if } d_\text{ind} < 0.25\,\text{m} \label{eq:ind} \\
r_\text{grp} &= (d_\text{grp,thr} - d_\text{grp}) \cdot 10 \cdot \Delta t
  \quad \text{if } d_\text{grp} < 0.35\,\text{m} \label{eq:grp}
\end{align}
$r_\text{pot}$ is always active ($\approx{+}0.5$/step at $v_\text{pref}=1$\,m/s).
Terminal events override $r_t$: success $+10$, collision $-20$, group intrusion $-5$.

%===============================================================================
\section{Group Detection Evaluation}
\label{sec:detection}

To validate the GroupDetector + SlotAttention backbone \emph{independently of navigation},
we evaluate it on 100 held-out episodes from the realistic benchmark.
Ground-truth group membership is read from \texttt{human.group\_id};
predicted pair scores are $\hat{y}_{ij} = \sum_k \alpha_{ki}\alpha_{kj}$.

\begin{table}[h]
\centering
\caption{\textbf{Group detection performance vs.\ DBSCAN baselines.}
DBSCAN uses position-only features; DBSCAN+V uses position and velocity.
Averaged over 100 episodes. \tbd{} = run \texttt{eval\_group\_detection.py} to fill.}
\label{tab:detection}
\small
\begin{tabular}{lccccc}
\toprule
\textbf{Method} & \textbf{Pair F1}$\uparrow$ & \textbf{AUC}$\uparrow$
  & \textbf{Purity}$\uparrow$ & \textbf{ARI}$\uparrow$ & \textbf{NMI}$\uparrow$ \\
\midrule
DBSCAN ($\epsilon$=0.5) & \tbd & \tbd & \tbd & \tbd & \tbd \\
DBSCAN ($\epsilon$=0.8) & \tbd & \tbd & \tbd & \tbd & \tbd \\
DBSCAN ($\epsilon$=1.0) & \tbd & \tbd & \tbd & \tbd & \tbd \\
DBSCAN ($\epsilon$=1.5) & \tbd & \tbd & \tbd & \tbd & \tbd \\
DBSCAN+V ($\epsilon$=0.8) & \tbd & \tbd & \tbd & \tbd & \tbd \\
\midrule
\best{\grace{} backbone (ours)} & \tbd & \tbd & \tbd & \tbd & \tbd \\
\bottomrule
\end{tabular}
\end{table}

\textit{Command:}
\texttt{python eval\_group\_detection.py -{}-model\_dir trained\_models/gram\_map/stageC
-{}-test\_model 41665.pt -{}-episodes 100 -{}-seed 0 -{}-out results/detection\_grace.csv}

\subsection{Detection Robustness Under Occlusion}
\label{sec:detection_occlusion}

We repeat the evaluation dropping 20\%, 40\%, and 60\% of visible humans at random
each step to simulate partial occlusion.

\begin{table}[h]
\centering
\caption{\textbf{Pair F1 under increasing occlusion.}
\grace{} is expected to degrade more gracefully than DBSCAN because the GNN
uses velocity context in addition to position.}
\label{tab:occlusion}
\small
\begin{tabular}{lccc}
\toprule
\textbf{Method} & \textbf{0\% dropped} & \textbf{40\% dropped} & \textbf{60\% dropped} \\
\midrule
DBSCAN (best $\epsilon$)  & \tbd & \tbd & \tbd \\
\grace{} backbone         & \tbd & \tbd & \tbd \\
\bottomrule
\end{tabular}
\end{table}

%===============================================================================
\section{Extended Ablation Analysis}
\label{sec:ablation_ext}

\subsection{Training Curves for All Ablation Variants}

\begin{figure}[h]
\centering
\figplaceholder{7cm}{%
  \textbf{Fig.~A\ref{fig:training_curves} — placeholder.}\\[4pt]
  Single plot: mean PPO reward vs.\ training updates.\\
  Lines: \grace{} full (black), C1 no-group (red), C3 no-traj (blue),
  C5 uniform-$\alpha$ (green), C4 no-aux seeds 0/1/2 (orange dashed).\\[4pt]
  \textit{Generate:} \texttt{python plot.py} — configure \texttt{model\_dirs}
  list to include all ablation output directories.\\
  Key visual: C4 orange lines should diverge widely, all others converge.%
}
\caption{\textbf{Training reward curves for \grace{} and all ablation variants.}
  C4 (no auxiliary loss) shows high variance across seeds; all other ablations
  converge, though at lower final reward than the full model.}
\label{fig:training_curves}
\end{figure}

\subsection{Per-Seed Breakdown for C4 (No Auxiliary Loss)}

\begin{table}[h]
\centering
\caption{\textbf{C4: per-seed navigation results.}
Three Stage~C training runs from the same Stage~B checkpoint, auxiliary loss disabled.
High variance confirms the instability claim in Section~6 of the main paper.}
\label{tab:c4_seeds}
\small
\begin{tabular}{lcccc}
\toprule
\textbf{Run} & \textbf{SR\%}$\uparrow$ & \textbf{CR\%}$\downarrow$
  & \textbf{GCR\%}$\downarrow$ & \textbf{Converged?} \\
\midrule
Seed 0              & \tbd & \tbd & \tbd & \tbd \\
Seed 1              & \tbd & \tbd & \tbd & \tbd \\
Seed 2              & \tbd & \tbd & \tbd & \tbd \\
\midrule
Mean $\pm$ std      & \tbd & \tbd & \tbd & — \\
\midrule
\grace{} full (ref) & $87\pm\tbd$ & $7\pm\tbd$ & $0\pm0$ & Yes \\
\bottomrule
\end{tabular}
\end{table}

\subsection{$K$-Slot Sweep}

\begin{figure}[h]
\centering
\figplaceholder{5cm}{%
  \textbf{Fig.~A\ref{fig:k_sweep} — placeholder.}\\[4pt]
  Dual y-axis line plot. x-axis: $K \in \{1,2,3,4,5\}$.
  Left y-axis: SR\% (inverted-U, peak at $K=3$).
  Right y-axis: GCR\% (U-shape, minimum at $K=3$).\\[4pt]
  \textit{Data:} \texttt{test.py} on
  \texttt{trained\_models/gram\_map/ablation\_K\{1,2,4,5\}} + full model ($K=3$).%
}
\caption{\textbf{Effect of slot count $K$ on SR and GCR.}
  $K=3$ achieves the best trade-off: higher SR than $K\leq2$ and lower GCR than
  $K=1$, with marginal degradation at $K>3$ due to over-segmentation.}
\label{fig:k_sweep}
\end{figure}

%===============================================================================
\section{Cost-Map Visualization Gallery}
\label{sec:gallery}

Fig.~\ref{fig:gallery} shows four critical frames from \texttt{visualize\_grace.py}.
Each row is one scenario: (left) ground-truth environment with group hulls;
(middle) robot's learned slot assignments; (right) all 9 cost-map channels.

\begin{figure}[h]
\centering
\figplaceholder{14cm}{%
  \textbf{Fig.~A\ref{fig:gallery} — 4-row gallery placeholder.}\\[6pt]
  Row 1 — Static F-formation blocking path.
    Expected: L6 sharp peak over centroid; robot curves around hull.\\[3pt]
  Row 2 — Tight corridor between two dynamic-LF groups.
    Expected: L7 halos nearly touching; planner threads the gap.\\[3pt]
  Row 3 — Mixed crowd (groups + individuals).
    Expected: middle panel shows 3 distinct slot colors for groups, grey
    for individuals.\\[3pt]
  Row 4 — Failure: bilateral flanking.
    Expected: L6+L7 saturate both sides; robot oscillates.\\[6pt]
  \textit{Generate:}
  \texttt{python visualize\_grace.py -{}-model\_dir trained\_models/gram\_map/stageC
    -{}-test\_model 41665.pt -{}-seed \{3,7,11,15\} -{}-max-steps 400}\\
  Extract key frame: \texttt{ffmpeg -i scene.mp4 -vframes 1 -ss \{T\} frame.png}%
}
\caption{\textbf{Cost-map gallery — four representative scenarios.}
  L6 (group cohesion) consistently locates group centroids;
  L7 (repulsion) provides the early-warning margin that routes the planner
  around group hulls without explicit programming.}
\label{fig:gallery}
\end{figure}

%===============================================================================
\section{Slot Assignment Analysis}
\label{sec:slots}

\subsection{Temporal Stability of Slot Assignments}

\begin{figure}[h]
\centering
\figplaceholder{5cm}{%
  \textbf{Fig.~A\ref{fig:slot_stability} — placeholder.}\\[4pt]
  x-axis: timestep (0–200). y-axis: mean assignment entropy
  $H = -\sum_k \alpha_{kn}\log\alpha_{kn}$ averaged over visible humans.\\
  Expected: high entropy ($\approx$0.9) in steps 0–8, stable below 0.3 afterwards.\\[4pt]
  \textit{Generate:} log \texttt{actor\_critic.base.\_last\_alpha} each step
  in \texttt{test.py}; compute entropy per human per step; plot mean $\pm$ std.%
}
\caption{\textbf{Slot assignment entropy over a full episode.}
  SlotAttention converges to stable, low-entropy group assignments within
  5--10 steps and maintains them for the episode, confirming temporal consistency.}
\label{fig:slot_stability}
\end{figure}

\subsection{Learned Assignments vs.\ Ground-Truth Groups}

\begin{figure}[h]
\centering
\figplaceholder{5cm}{%
  \textbf{Fig.~A\ref{fig:slot_gt} — placeholder.}\\[4pt]
  Side-by-side heatmaps at one timestep. Left: $\alpha \in \mathbb{R}^{3\times20}$
  (rows = slots, columns = humans, color = weight). Right: GT binary membership
  matrix (same layout, one bright column-block per group).\\
  Expected: each GT group block aligns with one dominant slot row on the left.\\[4pt]
  \textit{Extract:} \texttt{actor\_critic.base.\_last\_alpha[0].numpy()} in
  \texttt{visualize\_grace.py}; plot with \texttt{plt.imshow}.%
}
\caption{\textbf{Learned slot assignments vs.\ ground-truth group membership.}
  Each group is captured by one dominant slot; individuals receive distributed
  weights, confirming the competitive-softmax behaviour described in Section~3.2.}
\label{fig:slot_gt}
\end{figure}

%===============================================================================
\section{Failure Case Analysis}
\label{sec:failures}

\subsection{Failure Case 1 — Bilateral Group Corridor}

\begin{figure}[h]
\centering
\figplaceholder{5cm}{%
  \textbf{Fig.~A\ref{fig:failure_bilateral} — placeholder.}\\[4pt]
  Three-panel frame from \texttt{visualize\_grace.py} at the collision step.\\
  Scenario: two dynamic-LF groups converge from both sides simultaneously.\\
  Key: L6 (channel 5) shows two peaks bracketing the robot position;
  L7 (channel 6) halos overlap in the robot's forward direction.\\[4pt]
  \textit{Find seed:} sweep \texttt{-{}-seed} 20–40 until a bilateral-flank
  failure is observed in the left panel trajectory.%
}
\caption{\textbf{Failure case 1: bilateral group corridor.}
  Both L6 peaks are active with no navigable corridor between them.
  The robot oscillates and eventually grazes the nearer group hull.}
\label{fig:failure_bilateral}
\end{figure}

\subsection{Failure Case 2 — Dense Crowd Saturation}

\begin{figure}[h]
\centering
\figplaceholder{5cm}{%
  \textbf{Fig.~A\ref{fig:failure_saturation} — placeholder.}\\[4pt]
  Three-panel frame with \texttt{config.sim.human\_num = 28}.\\
  Key: L1 (individual occupancy) near-uniform ($>$0.8); L6 no longer the
  dominant channel — group cohesion signal is buried under individual noise.\\[4pt]
  \textit{Generate:} temporarily set \texttt{human\_num=28} in config,
  run \texttt{test.py}, capture a timeout episode frame.%
}
\caption{\textbf{Failure case 2: cost-map saturation at $N>25$ pedestrians.}
  Individual and trajectory channels saturate, masking the group cohesion layer.
  The planner reverts to individual-occupancy-based steering and times out.}
\label{fig:failure_saturation}
\end{figure}
 in grace.tex)
% Not a standalone document — no \documentclass or \begin{document}.
% ============================================================

% Reset counters so appendix figures/tables are numbered A1, A2, ...
\renewcommand{\thefigure}{A\arabic{figure}}
\renewcommand{\thetable}{A\arabic{table}}
\renewcommand{\theequation}{A\arabic{equation}}
\setcounter{figure}{0}
\setcounter{table}{0}
\setcounter{equation}{0}

% ── Placeholder helper (requires tcolorbox — loaded in grace.tex preamble) ────
\newcommand{\figplaceholder}[2]{%
  \begin{tcolorbox}[colback=gray!10,colframe=gray!50,
                    width=\linewidth,height=#1,
                    halign=center,valign=center]
    \textit{\small #2}
  \end{tcolorbox}%
}

\noindent
This appendix provides full implementation details~(\S\ref{sec:impl}),
group-detection evaluation against DBSCAN baselines~(\S\ref{sec:detection}),
extended ablation analysis with training curves and per-seed
breakdowns~(\S\ref{sec:ablation_ext}), a cost-map visualization
gallery~(\S\ref{sec:gallery}), slot assignment
analysis~(\S\ref{sec:slots}), and failure-case analysis~(\S\ref{sec:failures}).

%===============================================================================
\section{Implementation Details}
\label{sec:impl}

\subsection{Architecture Specifications}

Table~\ref{tab:arch} gives the full specification of every module in \grace{}.

\begin{table}[h]
\centering
\caption{Complete \grace{} architecture with parameter counts.}
\label{tab:arch}
\small
\begin{tabular}{llr}
\toprule
\textbf{Module} & \textbf{Specification} & \textbf{Params} \\
\midrule
\multicolumn{3}{l}{\textit{Perception backbone (pretrained, Stage A)}} \\
GroupDetector   & GNN, 3 layers, input 21-d, embed 64-d & \tbd \\
SlotAttention   & $K=3$ slots, 64-d, 3 iters            & \tbd \\
\midrule
\multicolumn{3}{l}{\textit{Cost-map synthesis (analytic, zero learnable params)}} \\
CostMapSynthesizer & 9 ch, $32{\times}32$ grid, 6\,m range & 0 \\
\midrule
\multicolumn{3}{l}{\textit{Planning head (trained, Stages B--C)}} \\
CostMapPlanner  & Conv 32→64→128→128, AdaptivePool(4)        & \tbd \\
Robot MLP       & Linear 9→128→64                            & \tbd \\
Fusion          & Linear 320→256, ReLU                       & \tbd \\
GRUCell         & 256→256                                    & \tbd \\
Actor head      & Linear 256→256, Tanh $\times$2             & \tbd \\
Critic head     & Linear 256→256, Tanh $\times$2, Linear 256→1 & \tbd \\
OccupancyHead   & Conv $C$→32→$T$ ($1{\times}1$), no sigmoid & \tbd \\
\midrule
\textbf{Total trainable (Stage C)} & & \tbd \\
\bottomrule
\end{tabular}
\end{table}

\subsection{Full Hyperparameter Table}

\begin{table}[h]
\centering
\caption{PPO and curriculum hyperparameters for all training stages.}
\label{tab:hparams}
\small
\begin{tabular}{llll}
\toprule
\textbf{Hyperparameter} & \textbf{Stage A} & \textbf{Stage B} & \textbf{Stage C} \\
\midrule
Algorithm           & Supervised          & PPO              & PPO \\
Learning rate       & $1\!\times\!10^{-3}$ & $5\!\times\!10^{-4}$ & $5\!\times\!10^{-5}$ \\
LR schedule         & Cosine              & Linear decay     & Linear decay \\
Parallel envs       & —                   & 16               & 16 \\
Rollout steps       & —                   & 30               & 30 \\
PPO epochs          & —                   & 5                & 5 \\
Mini-batches        & —                   & 2                & 2 \\
Clip $\epsilon$     & —                   & 0.2              & 0.2 \\
Discount $\gamma$   & —                   & 0.99             & 0.99 \\
GAE $\lambda$       & —                   & 0.95             & 0.95 \\
Entropy coef        & —                   & 0.00             & 0.00 \\
Value loss coef     & —                   & 0.5              & 0.5 \\
Grad clip norm      & —                   & 0.5              & 0.5 \\
Aux loss $\lambda$  & —                   & 0.0 (off)        & 0.1 \\
Backbone frozen     & —                   & Yes              & No  \\
Total env steps     & —                   & 20M              & 20M \\
\midrule
\multicolumn{4}{l}{\textit{Environment — Stage B}} \\
\# humans           & — & 10  & — \\
Circle radius (m)   & — & 5.0 & — \\
Groups              & — & 2×static-F & — \\
On-path groups      & — & 0   & — \\
Realistic phases    & — & Off & — \\
\midrule
\multicolumn{4}{l}{\textit{Environment — Stage C}} \\
\# humans           & — & — & 20 \\
Circle radius (m)   & — & — & 8.5 \\
Groups              & — & — & 3×mixed \\
On-path groups      & — & — & 2 \\
Realistic phases    & — & — & A–E (all on) \\
\bottomrule
\end{tabular}
\end{table}

\subsection{Reward Shaping Details}

The reward at each non-terminal step is $r_t = r_\text{pot} + r_\text{grp} + r_\text{ind}$, where:
\begin{align}
r_\text{pot} &= 2\,(\text{dist}_{t-1} - \text{dist}_t) \label{eq:pot} \\
r_\text{ind} &= (d_\text{thr} - d_\text{ind}) \cdot 10 \cdot \Delta t
  \quad \text{if } d_\text{ind} < 0.25\,\text{m} \label{eq:ind} \\
r_\text{grp} &= (d_\text{grp,thr} - d_\text{grp}) \cdot 10 \cdot \Delta t
  \quad \text{if } d_\text{grp} < 0.35\,\text{m} \label{eq:grp}
\end{align}
$r_\text{pot}$ is always active ($\approx{+}0.5$/step at $v_\text{pref}=1$\,m/s).
Terminal events override $r_t$: success $+10$, collision $-20$, group intrusion $-5$.

%===============================================================================
\section{Group Detection Evaluation}
\label{sec:detection}

To validate the GroupDetector + SlotAttention backbone \emph{independently of navigation},
we evaluate it on 100 held-out episodes from the realistic benchmark.
Ground-truth group membership is read from \texttt{human.group\_id};
predicted pair scores are $\hat{y}_{ij} = \sum_k \alpha_{ki}\alpha_{kj}$.

\begin{table}[h]
\centering
\caption{\textbf{Group detection performance vs.\ DBSCAN baselines.}
DBSCAN uses position-only features; DBSCAN+V uses position and velocity.
Averaged over 100 episodes. \tbd{} = run \texttt{eval\_group\_detection.py} to fill.}
\label{tab:detection}
\small
\begin{tabular}{lccccc}
\toprule
\textbf{Method} & \textbf{Pair F1}$\uparrow$ & \textbf{AUC}$\uparrow$
  & \textbf{Purity}$\uparrow$ & \textbf{ARI}$\uparrow$ & \textbf{NMI}$\uparrow$ \\
\midrule
DBSCAN ($\epsilon$=0.5) & \tbd & \tbd & \tbd & \tbd & \tbd \\
DBSCAN ($\epsilon$=0.8) & \tbd & \tbd & \tbd & \tbd & \tbd \\
DBSCAN ($\epsilon$=1.0) & \tbd & \tbd & \tbd & \tbd & \tbd \\
DBSCAN ($\epsilon$=1.5) & \tbd & \tbd & \tbd & \tbd & \tbd \\
DBSCAN+V ($\epsilon$=0.8) & \tbd & \tbd & \tbd & \tbd & \tbd \\
\midrule
\best{\grace{} backbone (ours)} & \tbd & \tbd & \tbd & \tbd & \tbd \\
\bottomrule
\end{tabular}
\end{table}

\textit{Command:}
\texttt{python eval\_group\_detection.py -{}-model\_dir trained\_models/gram\_map/stageC
-{}-test\_model 41665.pt -{}-episodes 100 -{}-seed 0 -{}-out results/detection\_grace.csv}

\subsection{Detection Robustness Under Occlusion}
\label{sec:detection_occlusion}

We repeat the evaluation dropping 20\%, 40\%, and 60\% of visible humans at random
each step to simulate partial occlusion.

\begin{table}[h]
\centering
\caption{\textbf{Pair F1 under increasing occlusion.}
\grace{} is expected to degrade more gracefully than DBSCAN because the GNN
uses velocity context in addition to position.}
\label{tab:occlusion}
\small
\begin{tabular}{lccc}
\toprule
\textbf{Method} & \textbf{0\% dropped} & \textbf{40\% dropped} & \textbf{60\% dropped} \\
\midrule
DBSCAN (best $\epsilon$)  & \tbd & \tbd & \tbd \\
\grace{} backbone         & \tbd & \tbd & \tbd \\
\bottomrule
\end{tabular}
\end{table}

%===============================================================================
\section{Extended Ablation Analysis}
\label{sec:ablation_ext}

\subsection{Training Curves for All Ablation Variants}

\begin{figure}[h]
\centering
\figplaceholder{7cm}{%
  \textbf{Fig.~A\ref{fig:training_curves} — placeholder.}\\[4pt]
  Single plot: mean PPO reward vs.\ training updates.\\
  Lines: \grace{} full (black), C1 no-group (red), C3 no-traj (blue),
  C5 uniform-$\alpha$ (green), C4 no-aux seeds 0/1/2 (orange dashed).\\[4pt]
  \textit{Generate:} \texttt{python plot.py} — configure \texttt{model\_dirs}
  list to include all ablation output directories.\\
  Key visual: C4 orange lines should diverge widely, all others converge.%
}
\caption{\textbf{Training reward curves for \grace{} and all ablation variants.}
  C4 (no auxiliary loss) shows high variance across seeds; all other ablations
  converge, though at lower final reward than the full model.}
\label{fig:training_curves}
\end{figure}

\subsection{Per-Seed Breakdown for C4 (No Auxiliary Loss)}

\begin{table}[h]
\centering
\caption{\textbf{C4: per-seed navigation results.}
Three Stage~C training runs from the same Stage~B checkpoint, auxiliary loss disabled.
High variance confirms the instability claim in Section~6 of the main paper.}
\label{tab:c4_seeds}
\small
\begin{tabular}{lcccc}
\toprule
\textbf{Run} & \textbf{SR\%}$\uparrow$ & \textbf{CR\%}$\downarrow$
  & \textbf{GCR\%}$\downarrow$ & \textbf{Converged?} \\
\midrule
Seed 0              & \tbd & \tbd & \tbd & \tbd \\
Seed 1              & \tbd & \tbd & \tbd & \tbd \\
Seed 2              & \tbd & \tbd & \tbd & \tbd \\
\midrule
Mean $\pm$ std      & \tbd & \tbd & \tbd & — \\
\midrule
\grace{} full (ref) & $87\pm\tbd$ & $7\pm\tbd$ & $0\pm0$ & Yes \\
\bottomrule
\end{tabular}
\end{table}

\subsection{$K$-Slot Sweep}

\begin{figure}[h]
\centering
\figplaceholder{5cm}{%
  \textbf{Fig.~A\ref{fig:k_sweep} — placeholder.}\\[4pt]
  Dual y-axis line plot. x-axis: $K \in \{1,2,3,4,5\}$.
  Left y-axis: SR\% (inverted-U, peak at $K=3$).
  Right y-axis: GCR\% (U-shape, minimum at $K=3$).\\[4pt]
  \textit{Data:} \texttt{test.py} on
  \texttt{trained\_models/gram\_map/ablation\_K\{1,2,4,5\}} + full model ($K=3$).%
}
\caption{\textbf{Effect of slot count $K$ on SR and GCR.}
  $K=3$ achieves the best trade-off: higher SR than $K\leq2$ and lower GCR than
  $K=1$, with marginal degradation at $K>3$ due to over-segmentation.}
\label{fig:k_sweep}
\end{figure}

%===============================================================================
\section{Cost-Map Visualization Gallery}
\label{sec:gallery}

Fig.~\ref{fig:gallery} shows four critical frames from \texttt{visualize\_grace.py}.
Each row is one scenario: (left) ground-truth environment with group hulls;
(middle) robot's learned slot assignments; (right) all 9 cost-map channels.

\begin{figure}[h]
\centering
\figplaceholder{14cm}{%
  \textbf{Fig.~A\ref{fig:gallery} — 4-row gallery placeholder.}\\[6pt]
  Row 1 — Static F-formation blocking path.
    Expected: L6 sharp peak over centroid; robot curves around hull.\\[3pt]
  Row 2 — Tight corridor between two dynamic-LF groups.
    Expected: L7 halos nearly touching; planner threads the gap.\\[3pt]
  Row 3 — Mixed crowd (groups + individuals).
    Expected: middle panel shows 3 distinct slot colors for groups, grey
    for individuals.\\[3pt]
  Row 4 — Failure: bilateral flanking.
    Expected: L6+L7 saturate both sides; robot oscillates.\\[6pt]
  \textit{Generate:}
  \texttt{python visualize\_grace.py -{}-model\_dir trained\_models/gram\_map/stageC
    -{}-test\_model 41665.pt -{}-seed \{3,7,11,15\} -{}-max-steps 400}\\
  Extract key frame: \texttt{ffmpeg -i scene.mp4 -vframes 1 -ss \{T\} frame.png}%
}
\caption{\textbf{Cost-map gallery — four representative scenarios.}
  L6 (group cohesion) consistently locates group centroids;
  L7 (repulsion) provides the early-warning margin that routes the planner
  around group hulls without explicit programming.}
\label{fig:gallery}
\end{figure}

%===============================================================================
\section{Slot Assignment Analysis}
\label{sec:slots}

\subsection{Temporal Stability of Slot Assignments}

\begin{figure}[h]
\centering
\figplaceholder{5cm}{%
  \textbf{Fig.~A\ref{fig:slot_stability} — placeholder.}\\[4pt]
  x-axis: timestep (0–200). y-axis: mean assignment entropy
  $H = -\sum_k \alpha_{kn}\log\alpha_{kn}$ averaged over visible humans.\\
  Expected: high entropy ($\approx$0.9) in steps 0–8, stable below 0.3 afterwards.\\[4pt]
  \textit{Generate:} log \texttt{actor\_critic.base.\_last\_alpha} each step
  in \texttt{test.py}; compute entropy per human per step; plot mean $\pm$ std.%
}
\caption{\textbf{Slot assignment entropy over a full episode.}
  SlotAttention converges to stable, low-entropy group assignments within
  5--10 steps and maintains them for the episode, confirming temporal consistency.}
\label{fig:slot_stability}
\end{figure}

\subsection{Learned Assignments vs.\ Ground-Truth Groups}

\begin{figure}[h]
\centering
\figplaceholder{5cm}{%
  \textbf{Fig.~A\ref{fig:slot_gt} — placeholder.}\\[4pt]
  Side-by-side heatmaps at one timestep. Left: $\alpha \in \mathbb{R}^{3\times20}$
  (rows = slots, columns = humans, color = weight). Right: GT binary membership
  matrix (same layout, one bright column-block per group).\\
  Expected: each GT group block aligns with one dominant slot row on the left.\\[4pt]
  \textit{Extract:} \texttt{actor\_critic.base.\_last\_alpha[0].numpy()} in
  \texttt{visualize\_grace.py}; plot with \texttt{plt.imshow}.%
}
\caption{\textbf{Learned slot assignments vs.\ ground-truth group membership.}
  Each group is captured by one dominant slot; individuals receive distributed
  weights, confirming the competitive-softmax behaviour described in Section~3.2.}
\label{fig:slot_gt}
\end{figure}

%===============================================================================
\section{Failure Case Analysis}
\label{sec:failures}

\subsection{Failure Case 1 — Bilateral Group Corridor}

\begin{figure}[h]
\centering
\figplaceholder{5cm}{%
  \textbf{Fig.~A\ref{fig:failure_bilateral} — placeholder.}\\[4pt]
  Three-panel frame from \texttt{visualize\_grace.py} at the collision step.\\
  Scenario: two dynamic-LF groups converge from both sides simultaneously.\\
  Key: L6 (channel 5) shows two peaks bracketing the robot position;
  L7 (channel 6) halos overlap in the robot's forward direction.\\[4pt]
  \textit{Find seed:} sweep \texttt{-{}-seed} 20–40 until a bilateral-flank
  failure is observed in the left panel trajectory.%
}
\caption{\textbf{Failure case 1: bilateral group corridor.}
  Both L6 peaks are active with no navigable corridor between them.
  The robot oscillates and eventually grazes the nearer group hull.}
\label{fig:failure_bilateral}
\end{figure}

\subsection{Failure Case 2 — Dense Crowd Saturation}

\begin{figure}[h]
\centering
\figplaceholder{5cm}{%
  \textbf{Fig.~A\ref{fig:failure_saturation} — placeholder.}\\[4pt]
  Three-panel frame with \texttt{config.sim.human\_num = 28}.\\
  Key: L1 (individual occupancy) near-uniform ($>$0.8); L6 no longer the
  dominant channel — group cohesion signal is buried under individual noise.\\[4pt]
  \textit{Generate:} temporarily set \texttt{human\_num=28} in config,
  run \texttt{test.py}, capture a timeout episode frame.%
}
\caption{\textbf{Failure case 2: cost-map saturation at $N>25$ pedestrians.}
  Individual and trajectory channels saturate, masking the group cohesion layer.
  The planner reverts to individual-occupancy-based steering and times out.}
\label{fig:failure_saturation}
\end{figure}
 in grace.tex)
% Not a standalone document — no \documentclass or \begin{document}.
% ============================================================

% Reset counters so appendix figures/tables are numbered A1, A2, ...
\renewcommand{\thefigure}{A\arabic{figure}}
\renewcommand{\thetable}{A\arabic{table}}
\renewcommand{\theequation}{A\arabic{equation}}
\setcounter{figure}{0}
\setcounter{table}{0}
\setcounter{equation}{0}

% ── Placeholder helper (requires tcolorbox — loaded in grace.tex preamble) ────
\newcommand{\figplaceholder}[2]{%
  \begin{tcolorbox}[colback=gray!10,colframe=gray!50,
                    width=\linewidth,height=#1,
                    halign=center,valign=center]
    \textit{\small #2}
  \end{tcolorbox}%
}

\noindent
This appendix provides full implementation details~(\S\ref{sec:impl}),
group-detection evaluation against DBSCAN baselines~(\S\ref{sec:detection}),
extended ablation analysis with training curves and per-seed
breakdowns~(\S\ref{sec:ablation_ext}), a cost-map visualization
gallery~(\S\ref{sec:gallery}), slot assignment
analysis~(\S\ref{sec:slots}), and failure-case analysis~(\S\ref{sec:failures}).

%===============================================================================
\section{Implementation Details}
\label{sec:impl}

\subsection{Architecture Specifications}

Table~\ref{tab:arch} gives the full specification of every module in \grace{}.

\begin{table}[h]
\centering
\caption{Complete \grace{} architecture with parameter counts.}
\label{tab:arch}
\small
\begin{tabular}{llr}
\toprule
\textbf{Module} & \textbf{Specification} & \textbf{Params} \\
\midrule
\multicolumn{3}{l}{\textit{Perception backbone (pretrained, Stage A)}} \\
GroupDetector   & GNN, 3 layers, input 21-d, embed 64-d & \tbd \\
SlotAttention   & $K=3$ slots, 64-d, 3 iters            & \tbd \\
\midrule
\multicolumn{3}{l}{\textit{Cost-map synthesis (analytic, zero learnable params)}} \\
CostMapSynthesizer & 9 ch, $32{\times}32$ grid, 6\,m range & 0 \\
\midrule
\multicolumn{3}{l}{\textit{Planning head (trained, Stages B--C)}} \\
CostMapPlanner  & Conv 32→64→128→128, AdaptivePool(4)        & \tbd \\
Robot MLP       & Linear 9→128→64                            & \tbd \\
Fusion          & Linear 320→256, ReLU                       & \tbd \\
GRUCell         & 256→256                                    & \tbd \\
Actor head      & Linear 256→256, Tanh $\times$2             & \tbd \\
Critic head     & Linear 256→256, Tanh $\times$2, Linear 256→1 & \tbd \\
OccupancyHead   & Conv $C$→32→$T$ ($1{\times}1$), no sigmoid & \tbd \\
\midrule
\textbf{Total trainable (Stage C)} & & \tbd \\
\bottomrule
\end{tabular}
\end{table}

\subsection{Full Hyperparameter Table}

\begin{table}[h]
\centering
\caption{PPO and curriculum hyperparameters for all training stages.}
\label{tab:hparams}
\small
\begin{tabular}{llll}
\toprule
\textbf{Hyperparameter} & \textbf{Stage A} & \textbf{Stage B} & \textbf{Stage C} \\
\midrule
Algorithm           & Supervised          & PPO              & PPO \\
Learning rate       & $1\!\times\!10^{-3}$ & $5\!\times\!10^{-4}$ & $5\!\times\!10^{-5}$ \\
LR schedule         & Cosine              & Linear decay     & Linear decay \\
Parallel envs       & —                   & 16               & 16 \\
Rollout steps       & —                   & 30               & 30 \\
PPO epochs          & —                   & 5                & 5 \\
Mini-batches        & —                   & 2                & 2 \\
Clip $\epsilon$     & —                   & 0.2              & 0.2 \\
Discount $\gamma$   & —                   & 0.99             & 0.99 \\
GAE $\lambda$       & —                   & 0.95             & 0.95 \\
Entropy coef        & —                   & 0.00             & 0.00 \\
Value loss coef     & —                   & 0.5              & 0.5 \\
Grad clip norm      & —                   & 0.5              & 0.5 \\
Aux loss $\lambda$  & —                   & 0.0 (off)        & 0.1 \\
Backbone frozen     & —                   & Yes              & No  \\
Total env steps     & —                   & 20M              & 20M \\
\midrule
\multicolumn{4}{l}{\textit{Environment — Stage B}} \\
\# humans           & — & 10  & — \\
Circle radius (m)   & — & 5.0 & — \\
Groups              & — & 2×static-F & — \\
On-path groups      & — & 0   & — \\
Realistic phases    & — & Off & — \\
\midrule
\multicolumn{4}{l}{\textit{Environment — Stage C}} \\
\# humans           & — & — & 20 \\
Circle radius (m)   & — & — & 8.5 \\
Groups              & — & — & 3×mixed \\
On-path groups      & — & — & 2 \\
Realistic phases    & — & — & A–E (all on) \\
\bottomrule
\end{tabular}
\end{table}

\subsection{Reward Shaping Details}

The reward at each non-terminal step is $r_t = r_\text{pot} + r_\text{grp} + r_\text{ind}$, where:
\begin{align}
r_\text{pot} &= 2\,(\text{dist}_{t-1} - \text{dist}_t) \label{eq:pot} \\
r_\text{ind} &= (d_\text{thr} - d_\text{ind}) \cdot 10 \cdot \Delta t
  \quad \text{if } d_\text{ind} < 0.25\,\text{m} \label{eq:ind} \\
r_\text{grp} &= (d_\text{grp,thr} - d_\text{grp}) \cdot 10 \cdot \Delta t
  \quad \text{if } d_\text{grp} < 0.35\,\text{m} \label{eq:grp}
\end{align}
$r_\text{pot}$ is always active ($\approx{+}0.5$/step at $v_\text{pref}=1$\,m/s).
Terminal events override $r_t$: success $+10$, collision $-20$, group intrusion $-5$.

%===============================================================================
\section{Group Detection Evaluation}
\label{sec:detection}

To validate the GroupDetector + SlotAttention backbone \emph{independently of navigation},
we evaluate it on 100 held-out episodes from the realistic benchmark.
Ground-truth group membership is read from \texttt{human.group\_id};
predicted pair scores are $\hat{y}_{ij} = \sum_k \alpha_{ki}\alpha_{kj}$.

\begin{table}[h]
\centering
\caption{\textbf{Group detection performance vs.\ DBSCAN baselines.}
DBSCAN uses position-only features; DBSCAN+V uses position and velocity.
Averaged over 100 episodes. \tbd{} = run \texttt{eval\_group\_detection.py} to fill.}
\label{tab:detection}
\small
\begin{tabular}{lccccc}
\toprule
\textbf{Method} & \textbf{Pair F1}$\uparrow$ & \textbf{AUC}$\uparrow$
  & \textbf{Purity}$\uparrow$ & \textbf{ARI}$\uparrow$ & \textbf{NMI}$\uparrow$ \\
\midrule
DBSCAN ($\epsilon$=0.5) & \tbd & \tbd & \tbd & \tbd & \tbd \\
DBSCAN ($\epsilon$=0.8) & \tbd & \tbd & \tbd & \tbd & \tbd \\
DBSCAN ($\epsilon$=1.0) & \tbd & \tbd & \tbd & \tbd & \tbd \\
DBSCAN ($\epsilon$=1.5) & \tbd & \tbd & \tbd & \tbd & \tbd \\
DBSCAN+V ($\epsilon$=0.8) & \tbd & \tbd & \tbd & \tbd & \tbd \\
\midrule
\best{\grace{} backbone (ours)} & \tbd & \tbd & \tbd & \tbd & \tbd \\
\bottomrule
\end{tabular}
\end{table}

\textit{Command:}
\texttt{python eval\_group\_detection.py -{}-model\_dir trained\_models/gram\_map/stageC
-{}-test\_model 41665.pt -{}-episodes 100 -{}-seed 0 -{}-out results/detection\_grace.csv}

\subsection{Detection Robustness Under Occlusion}
\label{sec:detection_occlusion}

We repeat the evaluation dropping 20\%, 40\%, and 60\% of visible humans at random
each step to simulate partial occlusion.

\begin{table}[h]
\centering
\caption{\textbf{Pair F1 under increasing occlusion.}
\grace{} is expected to degrade more gracefully than DBSCAN because the GNN
uses velocity context in addition to position.}
\label{tab:occlusion}
\small
\begin{tabular}{lccc}
\toprule
\textbf{Method} & \textbf{0\% dropped} & \textbf{40\% dropped} & \textbf{60\% dropped} \\
\midrule
DBSCAN (best $\epsilon$)  & \tbd & \tbd & \tbd \\
\grace{} backbone         & \tbd & \tbd & \tbd \\
\bottomrule
\end{tabular}
\end{table}

%===============================================================================
\section{Extended Ablation Analysis}
\label{sec:ablation_ext}

\subsection{Training Curves for All Ablation Variants}

\begin{figure}[h]
\centering
\figplaceholder{7cm}{%
  \textbf{Fig.~A\ref{fig:training_curves} — placeholder.}\\[4pt]
  Single plot: mean PPO reward vs.\ training updates.\\
  Lines: \grace{} full (black), C1 no-group (red), C3 no-traj (blue),
  C5 uniform-$\alpha$ (green), C4 no-aux seeds 0/1/2 (orange dashed).\\[4pt]
  \textit{Generate:} \texttt{python plot.py} — configure \texttt{model\_dirs}
  list to include all ablation output directories.\\
  Key visual: C4 orange lines should diverge widely, all others converge.%
}
\caption{\textbf{Training reward curves for \grace{} and all ablation variants.}
  C4 (no auxiliary loss) shows high variance across seeds; all other ablations
  converge, though at lower final reward than the full model.}
\label{fig:training_curves}
\end{figure}

\subsection{Per-Seed Breakdown for C4 (No Auxiliary Loss)}

\begin{table}[h]
\centering
\caption{\textbf{C4: per-seed navigation results.}
Three Stage~C training runs from the same Stage~B checkpoint, auxiliary loss disabled.
High variance confirms the instability claim in Section~6 of the main paper.}
\label{tab:c4_seeds}
\small
\begin{tabular}{lcccc}
\toprule
\textbf{Run} & \textbf{SR\%}$\uparrow$ & \textbf{CR\%}$\downarrow$
  & \textbf{GCR\%}$\downarrow$ & \textbf{Converged?} \\
\midrule
Seed 0              & \tbd & \tbd & \tbd & \tbd \\
Seed 1              & \tbd & \tbd & \tbd & \tbd \\
Seed 2              & \tbd & \tbd & \tbd & \tbd \\
\midrule
Mean $\pm$ std      & \tbd & \tbd & \tbd & — \\
\midrule
\grace{} full (ref) & $87\pm\tbd$ & $7\pm\tbd$ & $0\pm0$ & Yes \\
\bottomrule
\end{tabular}
\end{table}

\subsection{$K$-Slot Sweep}

\begin{figure}[h]
\centering
\figplaceholder{5cm}{%
  \textbf{Fig.~A\ref{fig:k_sweep} — placeholder.}\\[4pt]
  Dual y-axis line plot. x-axis: $K \in \{1,2,3,4,5\}$.
  Left y-axis: SR\% (inverted-U, peak at $K=3$).
  Right y-axis: GCR\% (U-shape, minimum at $K=3$).\\[4pt]
  \textit{Data:} \texttt{test.py} on
  \texttt{trained\_models/gram\_map/ablation\_K\{1,2,4,5\}} + full model ($K=3$).%
}
\caption{\textbf{Effect of slot count $K$ on SR and GCR.}
  $K=3$ achieves the best trade-off: higher SR than $K\leq2$ and lower GCR than
  $K=1$, with marginal degradation at $K>3$ due to over-segmentation.}
\label{fig:k_sweep}
\end{figure}

%===============================================================================
\section{Cost-Map Visualization Gallery}
\label{sec:gallery}

Fig.~\ref{fig:gallery} shows four critical frames from \texttt{visualize\_grace.py}.
Each row is one scenario: (left) ground-truth environment with group hulls;
(middle) robot's learned slot assignments; (right) all 9 cost-map channels.

\begin{figure}[h]
\centering
\figplaceholder{14cm}{%
  \textbf{Fig.~A\ref{fig:gallery} — 4-row gallery placeholder.}\\[6pt]
  Row 1 — Static F-formation blocking path.
    Expected: L6 sharp peak over centroid; robot curves around hull.\\[3pt]
  Row 2 — Tight corridor between two dynamic-LF groups.
    Expected: L7 halos nearly touching; planner threads the gap.\\[3pt]
  Row 3 — Mixed crowd (groups + individuals).
    Expected: middle panel shows 3 distinct slot colors for groups, grey
    for individuals.\\[3pt]
  Row 4 — Failure: bilateral flanking.
    Expected: L6+L7 saturate both sides; robot oscillates.\\[6pt]
  \textit{Generate:}
  \texttt{python visualize\_grace.py -{}-model\_dir trained\_models/gram\_map/stageC
    -{}-test\_model 41665.pt -{}-seed \{3,7,11,15\} -{}-max-steps 400}\\
  Extract key frame: \texttt{ffmpeg -i scene.mp4 -vframes 1 -ss \{T\} frame.png}%
}
\caption{\textbf{Cost-map gallery — four representative scenarios.}
  L6 (group cohesion) consistently locates group centroids;
  L7 (repulsion) provides the early-warning margin that routes the planner
  around group hulls without explicit programming.}
\label{fig:gallery}
\end{figure}

%===============================================================================
\section{Slot Assignment Analysis}
\label{sec:slots}

\subsection{Temporal Stability of Slot Assignments}

\begin{figure}[h]
\centering
\figplaceholder{5cm}{%
  \textbf{Fig.~A\ref{fig:slot_stability} — placeholder.}\\[4pt]
  x-axis: timestep (0–200). y-axis: mean assignment entropy
  $H = -\sum_k \alpha_{kn}\log\alpha_{kn}$ averaged over visible humans.\\
  Expected: high entropy ($\approx$0.9) in steps 0–8, stable below 0.3 afterwards.\\[4pt]
  \textit{Generate:} log \texttt{actor\_critic.base.\_last\_alpha} each step
  in \texttt{test.py}; compute entropy per human per step; plot mean $\pm$ std.%
}
\caption{\textbf{Slot assignment entropy over a full episode.}
  SlotAttention converges to stable, low-entropy group assignments within
  5--10 steps and maintains them for the episode, confirming temporal consistency.}
\label{fig:slot_stability}
\end{figure}

\subsection{Learned Assignments vs.\ Ground-Truth Groups}

\begin{figure}[h]
\centering
\figplaceholder{5cm}{%
  \textbf{Fig.~A\ref{fig:slot_gt} — placeholder.}\\[4pt]
  Side-by-side heatmaps at one timestep. Left: $\alpha \in \mathbb{R}^{3\times20}$
  (rows = slots, columns = humans, color = weight). Right: GT binary membership
  matrix (same layout, one bright column-block per group).\\
  Expected: each GT group block aligns with one dominant slot row on the left.\\[4pt]
  \textit{Extract:} \texttt{actor\_critic.base.\_last\_alpha[0].numpy()} in
  \texttt{visualize\_grace.py}; plot with \texttt{plt.imshow}.%
}
\caption{\textbf{Learned slot assignments vs.\ ground-truth group membership.}
  Each group is captured by one dominant slot; individuals receive distributed
  weights, confirming the competitive-softmax behaviour described in Section~3.2.}
\label{fig:slot_gt}
\end{figure}

%===============================================================================
\section{Failure Case Analysis}
\label{sec:failures}

\subsection{Failure Case 1 — Bilateral Group Corridor}

\begin{figure}[h]
\centering
\figplaceholder{5cm}{%
  \textbf{Fig.~A\ref{fig:failure_bilateral} — placeholder.}\\[4pt]
  Three-panel frame from \texttt{visualize\_grace.py} at the collision step.\\
  Scenario: two dynamic-LF groups converge from both sides simultaneously.\\
  Key: L6 (channel 5) shows two peaks bracketing the robot position;
  L7 (channel 6) halos overlap in the robot's forward direction.\\[4pt]
  \textit{Find seed:} sweep \texttt{-{}-seed} 20–40 until a bilateral-flank
  failure is observed in the left panel trajectory.%
}
\caption{\textbf{Failure case 1: bilateral group corridor.}
  Both L6 peaks are active with no navigable corridor between them.
  The robot oscillates and eventually grazes the nearer group hull.}
\label{fig:failure_bilateral}
\end{figure}

\subsection{Failure Case 2 — Dense Crowd Saturation}

\begin{figure}[h]
\centering
\figplaceholder{5cm}{%
  \textbf{Fig.~A\ref{fig:failure_saturation} — placeholder.}\\[4pt]
  Three-panel frame with \texttt{config.sim.human\_num = 28}.\\
  Key: L1 (individual occupancy) near-uniform ($>$0.8); L6 no longer the
  dominant channel — group cohesion signal is buried under individual noise.\\[4pt]
  \textit{Generate:} temporarily set \texttt{human\_num=28} in config,
  run \texttt{test.py}, capture a timeout episode frame.%
}
\caption{\textbf{Failure case 2: cost-map saturation at $N>25$ pedestrians.}
  Individual and trajectory channels saturate, masking the group cohesion layer.
  The planner reverts to individual-occupancy-based steering and times out.}
\label{fig:failure_saturation}
\end{figure}


\end{document}
